\section{System overview}
\label{sec:system}

This section describes the platform the method runs on: the arms and
their load-bearing geometry, the servo model, the software pipeline, the
pluggable IK layer, and the simulation-rehearsal harness used for both
practice and benchmarking.

\paragraph{Hardware.}
Two SO-101 arms are mounted \SI{0.30}{m} apart (bases at
\(y=\pm\SI{0.15}{m}\)) on a 3D-printed platform whose geometry is
generated from a single parametric design file; a MuJoCo/Isaac digital
twin is generated from the same source. Each arm has five revolute body
joints~--- shoulder pan, shoulder lift, elbow flex, wrist flex, wrist
roll~--- plus a gripper. Two geometric facts about this chain carry the
whole orientation story of this paper. First, the wrist-roll axis
coincides with the gripper's long (tip) axis, so rolling the gripper
about its own tip is a single joint's motion. Second, the lift, elbow,
and flex axes are mutually parallel: all three are horizontal axes
normal to the vertical plane the arm sweeps as the shoulder pan rotates
(the \emph{arm plane}). Rotations about parallel axes compose by adding
their angles, so the elevation of the gripper tip above the horizontal
is, up to a fixed offset, simply the signed sum of the lift, elbow, and
flex angles~--- an affine function of those three joints. The armplane
mapping (Section~\ref{sec:armplane}) relies on this to know that any
commanded tip elevation is reachable by the wrist-flex joint alone, and
the Telegrip wrist recovery (Section~\ref{sec:telegrip}) inverts the
same affine relation in closed form to set the wrist directly. The
shoulder-pan angle, by contrast, is the only joint that changes the
tip's azimuth, which is why end-effector yaw is coupled to base rotation
on this arm (Section~\ref{sec:prelim-retarget}).

\paragraph{Servo model.}
STS3215 serial bus servos are position-controlled; the digital twin
models them with position actuators (\(k_p=25\), \(k_v=1.5\)) plus joint
damping, armature, and friction matching the official SO-ARM100 MuJoCo
model.

\paragraph{Software pipeline.}
A reader thread streams Meta Quest controller poses (\SI{50}{Hz}),
buttons, grips, and triggers. Controller transforms are One-Euro
filtered (\(f_{c,\min}=0.8\), \(\beta=5.0\), \(f_{c,d}=0.9\);
Section~\ref{sec:prelim-filter}) in a shared data manager. The IK solver
thread (\SI{100}{Hz}) performs the retargeting of
Section~\ref{sec:method} and hands 10-DoF joint targets to per-arm joint
threads (real robot) or to the MuJoCo scene (simulation rehearsal).
Holding both grips engages teleoperation (the \emph{clutch},
Section~\ref{sec:clutch}); the first grip of a session calibrates the
handle axes; buttons toggle arm enablement and a home pose; triggers
drive the grippers. The trigger maps to a \emph{capped} jaw opening~---
fully pressed closes the jaws, fully released opens them to only a
fraction of their mechanical range~--- because the full splay is wider
than garment grasps need and makes fine pinching awkward; closing is
always complete, so grip strength is untouched.

\paragraph{Pluggable IK methods.}
The IK layer is a narrow interface: anchor to the measured joints,
accept per-frame end-effector pose targets, solve, and expose the
resulting configuration. The armplane pipeline's own tracker is a Pink
QP; a thin adapter exposes the same interface over any of the six
registered benchmark methods (Appendix~\ref{app:methods}) and adds a
joint-space rate limiter (\(\dot q_{\max}=\SI{3}{rad/s}\) in simulation,
\SI{2}{rad/s} on hardware). Every method can therefore be driven by the
real headset on the real arms, by the headset against the digital twin,
or by scripted mock trajectories in the benchmark~--- without touching
the real-time thread. Independently of this registry, the
\emph{unmodified upstream} Telegrip stack (web UI, WebXR streaming, its
own IK) can also drive the two arms directly for head-to-head
comparison, through a thin bridge that hands it this rig's serial ports
and motor calibrations.

\paragraph{Simulation rehearsal and benchmark harness.}
The rehearsal and deployment tools share \emph{one} control stack by
construction: they build the same filter, clutch, retargeting, envelope,
and IK layer from a single wiring routine and differ only in the robot
backend~--- the digital twin in rehearsal, the serial motor buses on the
real rig. A method therefore behaves identically in both, which is the
whole point of rehearsing before touching hardware. The rehearsal tool
runs that shared thread against the digital twin, with either the real
headset or a scripted mock device (patterns: table circles, wrist
oscillation, envelope excursion). By default it loads the complete twin
rig scene~--- printed board, adapter plates, camera tower, wrist
cameras~--- with \emph{all} collisions enabled (arm--arm, arm--rig,
arm--table; the nested base/shoulder mesh pair is explicitly excluded
since those parts interpenetrate by design). It renders the three
mounted cameras into live view windows and draws the world-frame triad
and the calibrated ``headset centre'' marker, so the operator can see
exactly where the system thinks the VR frame sits; it also draws a pair
of orientation triads per arm~--- one for the \emph{measured} gripper
pose and one for the \emph{commanded} target pose~--- so any attitude
error between what the arm does and what the mapping asks for is visible
at a glance. The
benchmark runner sweeps mock trajectories across methods and reports the
metrics of Section~\ref{sec:experiments}; a second runner does the same
for the out-of-envelope policy sweep. During this work we found and
fixed a rehearsal-harness artifact: while teleoperation was idle, the
harness fed measured (gravity-sagged) joints back as servo targets,
letting the arms ratchet \(\approx\)\SI{13}{cm} downward before
activation. Idle targets are now held constant, matching the real
robot's behaviour, which only writes servo goals while teleoperation is
active.
